\scanpage{28}%
\mspage{13}\space
\uline{sans fibre additive}, \uline{munies de formes différentielles invariantes} $\pi$ \uline{et} $\pi'$, \uline{et de même invariant modulaire}. \uline{On suppose que} $2$, $c_4$ \uline{et} $c_6$ \uline{sont inversibles}, \uline{i.e.} \uline{que} $2^{-1}$, $1/j$ et $1/\inserted{\handwritten{(}}j-1728\inserted{\handwritten{)}}$ \uline{sont des sections de} \mbox{\begingroup\edef\encoutersep{\the\dimexpr.5\pgflinewidth\relax}\setbox0=\hbox{$S$}\edef\encwda{\the\wd0}\edef\enchta{\the\ht0}\edef\encdpa{\the\dp0}\rlap{\smash{\tikz[baseline=0pt]{\begin{scope}[shift={(\the\dimexpr \encwda/2\relax,\the\dimexpr -\encdpa/2+\enchta/2\relax)}]\coordinate (enc-ess-center) at (0pt,0pt);\handwritten{\draw[thin,gray,rotate around={-20:(enc-ess-center)}] (enc-ess-center) ellipse (0.52em and 1.02em);}\end{scope}\pgfresetboundingbox\path[use as bounding box] (0pt,\the\dimexpr-\encdpa\relax) rectangle (\encwda,\enchta);}}}\copy0\endgroup}\reversemarginpar\margin{\inserted{\handwritten{\mbox{\begingroup\edef\encoutersep{\the\dimexpr.5\pgflinewidth\relax}\setbox0=\hbox{$\underline{O}_S$}\edef\encwda{\the\wd0}\edef\enchta{\the\ht0}\edef\encdpa{\the\dp0}\rlap{\smash{\tikz[baseline=0pt]{\begin{scope}[shift={(\the\dimexpr \encwda/2\relax,\the\dimexpr -\encdpa/2+\enchta/2\relax)}]\coordinate (enc-oh-center) at (0pt,0pt);\handwritten{\draw[thin,gray,rotate around={-15:(enc-oh-center)}] (enc-oh-center) ellipse (1.13em and 1.55em);}\end{scope}\pgfresetboundingbox\path[use as bounding box] (0pt,\the\dimexpr-\encdpa\relax) rectangle (\encwda,\enchta);}}}\copy0\endgroup}}}}\normalmarginpar. \uline{Alors}, \uline{le schéma} $\text{\uline{Isom}}(E,F)$ \uline{s'identifie au schéma défini par l'équation}
\[
\lambda^2=c_4\,c'_6/c_6\,c'_4 \, .
\]

En vertu de (3.6), se donner $(E,\pi)$ revient à se donner $a_2$, $a_4$ et $a_6$ vérifiant $a_2\,a_4-9\,a_6=0$, et un isomorphisme de $E$ sur $E'$ est décrit par $\lambda$ tel que $(E\,\pi)$ soit isomorphe à $(E',\lambda\,\pi')$, i.e. par $\lambda$ vérifiant
\[
a'_2=\lambda^2\,a_2 \qquad a'_4=\lambda^4\,a_4 \qquad a'_6=\lambda^6\,a_6 \, . \tag*{(6.4.1)}
\]

Par hypothèse, $j=j'$, de sorte que $(c_4^3,c_6^2,\Delta)$ est proportionnel à $(c_4^{\raisebox{-.38ex}{$\scriptstyle\prime$}\,3},c_6^{\raisebox{-.38ex}{$\scriptstyle\prime$}\,2},\Delta')$, et que $\mu=c_4\,c'_6/\raisebox{-.52ex}{$c_6\,c'_4$}$ vérifie
\[
c'_4=\mu^2\,c_4 \qquad c'_6=\mu^3\,c_6 \qquad \Delta'=\mu^6\,\Delta \, . \tag*{(6.4.2)}
\]

En vertu de (3.4), on a alors
\[
a'_2=\mu\,a_2 \qquad a'_4=\mu^2\,a_4 \qquad a'_6=\mu^3\,a_6 \, .
\]

De plus, d'après (3.4) encore, localement, soit $a_2$, soit $a_4$ et $a_6$ sont inversibles, de sorte que (6.4.1) équivaut à $\lambda^2=\mu$.

Pour prouver (III) en général, on peut se ramener à supposer $S$ noethérien, et supposer soit $2$ inversible (cas traité en (6.4)), soit%
% archive sheet number: 27
